% Frozen source-aware proof packet, 2026-08-25.
%
% Authenticated primary-source dependency:
% Sebastien Boucksom and Mattias Jonsson, Global pluripotential theory over
% a trivially valued field, Ann. Fac. Sci. Toulouse 31 (2022), 647--836.
% Local primary PDF: tmp/pdfs/bj22_global_trivial.pdf.
% SHA-256: 8d5975337f4ac88a475929ea6a436ec4a15d08715d5582bc0ba0eb3e07152c52.
% Used exactly: Definition 1.26, Lemma 4.46, equation (11.4),
% Proposition 11.6(ii), Lemma 11.7, and Question 11.16.
%
% Scope firewall:
% * The proof concerns the full weak and strong topologies of X^{lin}; it
%   does not contradict their coincidence on one fixed dual complex.
% * The varying valuations lie on blowups at pairwise distinct centers, so
%   no single fixed dual complex contains the tail of the sequence.
% * No resolution of singularities is used.

\begin{proof}
All the valuations in the statement have linear growth.  This follows from
the explicit computation of \(T_\omega\) below; it also follows for
\(v_j\) and \(v\) from the general fact that divisorial valuations have
linear growth.

We first prove weak convergence.  Let \(f\in k(X)^\times\), and put
\[
                         r=\operatorname{ord}_D(f)\in\mathbf Z.
\]
At the generic point of \(D\), divide \(f\) by the \(r\)-th power of a
local equation of \(D\).  The resulting element has a nonzero rational
residue
\[
                         \overline f_D\in k(D)^\times.
\]
The zero and pole sets of \(\overline f_D\) on the projective line \(D\)
are finite.  Away from their union, the divided element is a unit in
\(\mathcal O_{X,p}\), while a local equation of the smooth divisor \(D\)
has order one at \(p\).  Hence, for all sufficiently large \(j\),
\[
                  v_j(f)=\operatorname{ord}_{p_j}(f)
                        =r
                        =\operatorname{ord}_D(f)=v(f).
\tag{1}\label{eq:p11-weak-eventual}
\]
Here pairwise distinctness of the \(p_j\) guarantees that the sequence is
eventually outside every prescribed finite subset of \(D(k)\).  The weak
topology on the valuation locus is the topology of pointwise convergence
on \(k(X)^\times\).  Thus \eqref{eq:p11-weak-eventual} proves
\(v_j\to v\) weakly.

We next compute the exhaustion.  Fix \(j\).  Choose homogeneous linear
coordinates \(X_j,Y_j,Z_j\) such that
\[
            D=(X_j=0),\qquad p_j=(X_j=Y_j=0),
            \qquad Z_j(p_j)\ne0.
\]
For \(m\geq1\), every nonzero section
\(s\in H^0(X,\mathcal O_X(m))\), written in the local trivialization
\(Z_j^m\), is a nonzero polynomial
\[
       \frac{s}{Z_j^m}
          =\sum_{i+q\leq m}c_{iq}
             \left(\frac{X_j}{Z_j}\right)^i
             \left(\frac{Y_j}{Z_j}\right)^q.
\tag{2}\label{eq:p11-local-polynomial}
\]
Its order at \(p_j\) is the least \(i+q\) for which \(c_{iq}\ne0\), and
is therefore at most \(m\).  Equality is attained by \(s=X_j^m\).  By
Definition 1.26 and Lemma 4.46 of the cited source,
\[
                         T_\omega(v_j)=1.
\tag{3}\label{eq:p11-T-point}
\]
Likewise, the order along \(D\) of a degree-\(m\) homogeneous polynomial
is at most \(m\), with equality for \(X_j^m\).  Hence
\[
                         T_\omega(v)=1.
\tag{4}\label{eq:p11-T-divisor}
\]

Finally, let \(H_j=(Y_j=0)\).  This line contains \(p_j\) and is distinct
from \(D\).  For the section \(Y_j\in H^0(X,L)\) one has
\[
                   v_j(Y_j)=1,
                   \qquad v(Y_j)=\operatorname{ord}_D(Y_j)=0.
\tag{5}\label{eq:p11-moving-section}
\]
The section formula of Lemma 11.7 in the cited source gives
\[
\begin{aligned}
 d_{\infty,\omega}(v_j,v)
 &=\sup_{m\geq1}\ \sup_{0\ne s\in H^0(X,mL)}
       \frac{1}{m}\lvert v_j(s)-v(s)\rvert \\
 &\geq \lvert v_j(Y_j)-v(Y_j)\rvert=1.
\end{aligned}
\tag{6}\label{eq:p11-strong-gap}
\]
Thus the sequence is not strongly convergent although it is weakly
convergent and its exhaustion values agree identically with that of the
limit.  This is exactly a negative answer to the implication asked for in
Question 11.16 of the cited source.
\end{proof}
